\subsection*{iii.}

\addcontentsline{toc}{subsection}{iii.}

\textbf{¿Qué es una llave primaria?, ¿qué es una llave candidata?, ¿qué es una llave mínima?, ¿qué es una
super llave?}

\begin{itemize}
    \item \textbf{Llave primaria:} Es una llave candidata elegida como el medio principal para identificar tuplas dentro de una relación.
    \item \textbf{Llave candidata:} Es una super llave es un subconjunto del conjunto de atributos de la relación, a la cual no se le puede eliminar ningún atributo sin perder la capacidad de identificar de forma única a cada tupla.
    \item \textbf{Llave mínima:} Es una super llave para las cuales ningún subconjunto propio es una superllave.
    \item  \textbf{Super llave:} Conjunto de uno o más atributos cuyos valores no se repiten en tuplas distintas de la relación, los cuales permiten identificar de manera única una tupla en la relación.
\end{itemize}